% !TeX root = ../../main.tex
\section{循环坐标、泊松括号和正则变换}

\begin{solution}[泊松括号]
\problem{证明
\[
\{\omega,\lambda\}=-\{\lambda,\omega\},
\]
\[
\{\omega,\lambda+\sigma\}=\{\omega,\lambda\}+\{\omega,\sigma\},
\]
\[
\{\omega,\lambda\sigma\}=\{\omega,\lambda\}\sigma+\lambda\{\omega,\sigma\}.
\]
请注意上述公式与对易式 \((1.5.10)\) 和 \((1.5.11)\) 式之间的相似性。}
\answer
设 \(\omega,\lambda,\sigma\) 都是相空间变量 \(q,p\) 的函数。为简化记号，先在单自由度情形下证明；多自由度情形只需对每一对共轭变量 \((q_i,p_i)\) 求和，结论不变。单自由度泊松括号定义为
\[
\{\omega,\lambda\}
=
\frac{\partial \omega}{\partial q}\frac{\partial \lambda}{\partial p}
-
\frac{\partial \omega}{\partial p}\frac{\partial \lambda}{\partial q}.
\]
首先证明反对称性。由定义有
\[
\{\lambda,\omega\}
=
\frac{\partial \lambda}{\partial q}\frac{\partial \omega}{\partial p}
-
\frac{\partial \lambda}{\partial p}\frac{\partial \omega}{\partial q}.
\]
因此
\[
-\{\lambda,\omega\}
=
-\frac{\partial \lambda}{\partial q}\frac{\partial \omega}{\partial p}
+
\frac{\partial \lambda}{\partial p}\frac{\partial \omega}{\partial q}
=
\frac{\partial \omega}{\partial q}\frac{\partial \lambda}{\partial p}
-
\frac{\partial \omega}{\partial p}\frac{\partial \lambda}{\partial q}
=
\{\omega,\lambda\}.
\]
所以
\[
\{\omega,\lambda\}=-\{\lambda,\omega\}.
\]

其次证明对第二个变量的线性性。由定义，
\[
\{\omega,\lambda+\sigma\}
=
\frac{\partial \omega}{\partial q}\frac{\partial(\lambda+\sigma)}{\partial p}
-
\frac{\partial \omega}{\partial p}\frac{\partial(\lambda+\sigma)}{\partial q}.
\]
利用偏导数的线性性，
\[
\frac{\partial(\lambda+\sigma)}{\partial p}
=
\frac{\partial\lambda}{\partial p}
+
\frac{\partial\sigma}{\partial p},
\qquad
\frac{\partial(\lambda+\sigma)}{\partial q}
=
\frac{\partial\lambda}{\partial q}
+
\frac{\partial\sigma}{\partial q}.
\]
代入并整理可得
\[
\begin{aligned}
\{\omega,\lambda+\sigma\}
&=
\frac{\partial \omega}{\partial q}
\left(
\frac{\partial\lambda}{\partial p}
+
\frac{\partial\sigma}{\partial p}
\right)
-
\frac{\partial \omega}{\partial p}
\left(
\frac{\partial\lambda}{\partial q}
+
\frac{\partial\sigma}{\partial q}
\right)\\
&=
\left(
\frac{\partial \omega}{\partial q}\frac{\partial\lambda}{\partial p}
-
\frac{\partial \omega}{\partial p}\frac{\partial\lambda}{\partial q}
\right)
+
\left(
\frac{\partial \omega}{\partial q}\frac{\partial\sigma}{\partial p}
-
\frac{\partial \omega}{\partial p}\frac{\partial\sigma}{\partial q}
\right)\\
&=
\{\omega,\lambda\}+\{\omega,\sigma\}.
\end{aligned}
\]

最后证明乘积法则。由定义，
\[
\{\omega,\lambda\sigma\}
=
\frac{\partial \omega}{\partial q}\frac{\partial(\lambda\sigma)}{\partial p}
-
\frac{\partial \omega}{\partial p}\frac{\partial(\lambda\sigma)}{\partial q}.
\]
利用普通偏导数的乘积法则，
\[
\frac{\partial(\lambda\sigma)}{\partial p}
=
\frac{\partial\lambda}{\partial p}\sigma
+
\lambda\frac{\partial\sigma}{\partial p},
\qquad
\frac{\partial(\lambda\sigma)}{\partial q}
=
\frac{\partial\lambda}{\partial q}\sigma
+
\lambda\frac{\partial\sigma}{\partial q}.
\]
代入泊松括号定义，得到
\[
\begin{aligned}
\{\omega,\lambda\sigma\}
&=
\frac{\partial \omega}{\partial q}
\left(
\frac{\partial\lambda}{\partial p}\sigma
+
\lambda\frac{\partial\sigma}{\partial p}
\right)
-
\frac{\partial \omega}{\partial p}
\left(
\frac{\partial\lambda}{\partial q}\sigma
+
\lambda\frac{\partial\sigma}{\partial q}
\right)\\
&=
\left(
\frac{\partial \omega}{\partial q}\frac{\partial\lambda}{\partial p}
-
\frac{\partial \omega}{\partial p}\frac{\partial\lambda}{\partial q}
\right)\sigma
+
\lambda
\left(
\frac{\partial \omega}{\partial q}\frac{\partial\sigma}{\partial p}
-
\frac{\partial \omega}{\partial p}\frac{\partial\sigma}{\partial q}
\right)\\
&=
\{\omega,\lambda\}\sigma+\lambda\{\omega,\sigma\}.
\end{aligned}
\]
因此三条恒等式均得证。它们分别表明泊松括号具有反对称性、线性性和类似导数的乘积法则，这正是其与量子力学中对易式代数性质相似的原因。
\end{solution}

\begin{solution}[正则泊松括号关系]
\problem{(i) 证明 \((2.7.4)\) 和 \((2.7.5)\) 式。(ii) 考虑一个由 \(\mathcal H=p_x^2+p_y^2+ax^2+by^2\) 给出的二维问题。论证如果 \(a=b\)，则 \(\{l_z,\mathcal H\}\) 一定为零。用显式的计算证明。}
\answer
先证明正则变量的基本泊松括号关系。对 \(n\) 个自由度的系统，泊松括号定义为
\[
\{f,g\}=\sum_k\left(\frac{\partial f}{\partial q_k}\frac{\partial g}{\partial p_k}-\frac{\partial f}{\partial p_k}\frac{\partial g}{\partial q_k}\right).
\]
取 \(f=q_i,\ g=q_j\)，由于
\[
\frac{\partial q_i}{\partial q_k}=\delta_{ik},\qquad
\frac{\partial q_j}{\partial p_k}=0,\qquad
\frac{\partial q_i}{\partial p_k}=0,\qquad
\frac{\partial q_j}{\partial q_k}=\delta_{jk},
\]
所以
\[
\{q_i,q_j\}
=\sum_k(\delta_{ik}\cdot 0-0\cdot\delta_{jk})=0.
\]
同理，取 \(f=p_i,\ g=p_j\)，由于
\[
\frac{\partial p_i}{\partial q_k}=0,\qquad
\frac{\partial p_j}{\partial p_k}=\delta_{jk},\qquad
\frac{\partial p_i}{\partial p_k}=\delta_{ik},\qquad
\frac{\partial p_j}{\partial q_k}=0,
\]
所以
\[
\{p_i,p_j\}
=\sum_k(0\cdot\delta_{jk}-\delta_{ik}\cdot 0)=0.
\]
再取 \(f=q_i,\ g=p_j\)，得到
\[
\{q_i,p_j\}
=
\sum_k\left(
\frac{\partial q_i}{\partial q_k}\frac{\partial p_j}{\partial p_k}
-
\frac{\partial q_i}{\partial p_k}\frac{\partial p_j}{\partial q_k}
\right)
=
\sum_k(\delta_{ik}\delta_{jk}-0)
=
\delta_{ij}.
\]
因此基本正则泊松括号关系为
\[
\{q_i,q_j\}=0,\qquad
\{p_i,p_j\}=0,\qquad
\{q_i,p_j\}=\delta_{ij}.
\]
由反对称性还可得
\[
\{p_i,q_j\}=-\delta_{ij}.
\]
接着证明泊松括号与哈密顿正则方程的一致性。由定义，
\[
\{q_i,\mathcal H\}
=
\sum_k\left(
\frac{\partial q_i}{\partial q_k}\frac{\partial\mathcal H}{\partial p_k}
-
\frac{\partial q_i}{\partial p_k}\frac{\partial\mathcal H}{\partial q_k}
\right)
=
\frac{\partial \mathcal H}{\partial p_i}.
\]
根据哈密顿方程 \(\dot q_i=\partial\mathcal H/\partial p_i\)，可得
\[
\dot q_i=\{q_i,\mathcal H\}.
\]
类似地，
\[
\{p_i,\mathcal H\}
=
\sum_k\left(
\frac{\partial p_i}{\partial q_k}\frac{\partial\mathcal H}{\partial p_k}
-
\frac{\partial p_i}{\partial p_k}\frac{\partial\mathcal H}{\partial q_k}
\right)
=
-\frac{\partial \mathcal H}{\partial q_i}.
\]
根据哈密顿方程 \(\dot p_i=-\partial\mathcal H/\partial q_i\)，可得
\[
\dot p_i=\{p_i,\mathcal H\}.
\]
下面考虑二维哈密顿量
\[
\mathcal H=p_x^2+p_y^2+ax^2+by^2.
\]
绕 \(z\) 轴的角动量为
\[
l_z=xp_y-yp_x.
\]
如果 \(a=b\)，则哈密顿量变为
\[
\mathcal H=p_x^2+p_y^2+a(x^2+y^2).
\]
此时 \(\mathcal H\) 只依赖于 \(x^2+y^2\) 和 \(p_x^2+p_y^2\)，在 \(x\)-\(y\) 平面内具有旋转对称性。由于 \(l_z\) 是平面转动的生成元，旋转对称性意味着 \(l_z\) 守恒，因此应有
\[
\{l_z,\mathcal H\}=0.
\]
下面直接计算验证。按二维泊松括号定义，
\[
\{l_z,\mathcal H\}
=
\frac{\partial l_z}{\partial x}\frac{\partial\mathcal H}{\partial p_x}
+
\frac{\partial l_z}{\partial y}\frac{\partial\mathcal H}{\partial p_y}
-
\frac{\partial l_z}{\partial p_x}\frac{\partial\mathcal H}{\partial x}
-
\frac{\partial l_z}{\partial p_y}\frac{\partial\mathcal H}{\partial y}.
\]
由
\[
\frac{\partial l_z}{\partial x}=p_y,\qquad
\frac{\partial l_z}{\partial y}=-p_x,\qquad
\frac{\partial l_z}{\partial p_x}=-y,\qquad
\frac{\partial l_z}{\partial p_y}=x,
\]
以及
\[
\frac{\partial\mathcal H}{\partial p_x}=2p_x,\qquad
\frac{\partial\mathcal H}{\partial p_y}=2p_y,\qquad
\frac{\partial\mathcal H}{\partial x}=2ax,\qquad
\frac{\partial\mathcal H}{\partial y}=2by,
\]
可得
\[
\begin{aligned}
\{l_z,\mathcal H\}
&=
p_y(2p_x)+(-p_x)(2p_y)-(-y)(2ax)-x(2by)\\
&=
2p_xp_y-2p_xp_y+2axy-2bxy\\
&=
2(a-b)xy.
\end{aligned}
\]
因此当 \(a=b\) 时，
\[
\{l_z,\mathcal H\}=0.
\]
这说明在各向同性情形 \(a=b\) 下，系统具有绕 \(z\) 轴的旋转对称性，与该对称性对应的守恒量正是 \(l_z\)。
\end{solution}


\begin{solution}[正则泊松括号证明]
\problem{填补由 \((2.7.14)\) 式导出 \((2.7.18)\) 式所丢失的步骤。}
\answer
设 \((q_i,p_i)\) 是原来的正则坐标和正则动量，\((\bar q_i,\bar p_i)\) 是经过正则变换后得到的新变量。泊松括号在原变量下定义为
\[
\{f,g\}
=
\sum_i
\left(
\frac{\partial f}{\partial q_i}\frac{\partial g}{\partial p_i}
-
\frac{\partial f}{\partial p_i}\frac{\partial g}{\partial q_i}
\right).
\]
由 \((2.7.14)\) 式可知，新坐标 \(\bar q_j\) 的时间演化可写为
\[
\dot{\bar q}_j
=
\sum_i
\left(
\frac{\partial \bar q_j}{\partial q_i}\frac{\partial \mathcal H}{\partial p_i}
-
\frac{\partial \bar q_j}{\partial p_i}\frac{\partial \mathcal H}{\partial q_i}
\right).
\]
根据泊松括号定义，右边正是
\[
\dot{\bar q}_j=\{\bar q_j,\mathcal H\}.
\]
另一方面，将哈密顿量看成新变量的函数
\[
\mathcal H=\mathcal H(\bar q_1,\cdots,\bar q_n;\bar p_1,\cdots,\bar p_n),
\]
则由泊松括号的链式法则可得
\[
\{\bar q_j,\mathcal H\}
=
\sum_k
\left(
\frac{\partial \mathcal H}{\partial \bar q_k}
\{\bar q_j,\bar q_k\}
+
\frac{\partial \mathcal H}{\partial \bar p_k}
\{\bar q_j,\bar p_k\}
\right).
\]
如果变换后仍满足哈密顿正则方程，则必须有
\[
\dot{\bar q}_j=\frac{\partial \mathcal H}{\partial \bar p_j}.
\]
因此
\[
\sum_k
\left(
\frac{\partial \mathcal H}{\partial \bar q_k}
\{\bar q_j,\bar q_k\}
+
\frac{\partial \mathcal H}{\partial \bar p_k}
\{\bar q_j,\bar p_k\}
\right)
=
\frac{\partial \mathcal H}{\partial \bar p_j}.
\]
由于这个等式应对任意哈密顿量 \(\mathcal H\) 成立，\(\partial \mathcal H/\partial \bar q_k\) 和 \(\partial \mathcal H/\partial \bar p_k\) 的系数必须分别相等，于是得到
\[
\{\bar q_j,\bar q_k\}=0,\qquad
\{\bar q_j,\bar p_k\}=\delta_{jk}.
\]
同理，从新动量的时间演化出发。由 \((2.7.17)\) 式可写成
\[
\dot{\bar p}_j=\{\bar p_j,\mathcal H\}.
\]
再对 \(\mathcal H(\bar q,\bar p)\) 使用链式法则，有
\[
\{\bar p_j,\mathcal H\}
=
\sum_k
\left(
\frac{\partial \mathcal H}{\partial \bar q_k}
\{\bar p_j,\bar q_k\}
+
\frac{\partial \mathcal H}{\partial \bar p_k}
\{\bar p_j,\bar p_k\}
\right).
\]
若新变量仍满足哈密顿正则方程，则
\[
\dot{\bar p}_j=-\frac{\partial \mathcal H}{\partial \bar q_j}.
\]
因此
\[
\sum_k
\left(
\frac{\partial \mathcal H}{\partial \bar q_k}
\{\bar p_j,\bar q_k\}
+
\frac{\partial \mathcal H}{\partial \bar p_k}
\{\bar p_j,\bar p_k\}
\right)
=
-\frac{\partial \mathcal H}{\partial \bar q_j}.
\]
同样，由于该式对任意 \(\mathcal H\) 成立，比较各偏导数前的系数，得到
\[
\{\bar p_j,\bar q_k\}=-\delta_{jk},\qquad
\{\bar p_j,\bar p_k\}=0.
\]
综上，新变量满足
\[
\{\bar q_j,\bar q_k\}=0,\qquad
\{\bar p_j,\bar p_k\}=0,\qquad
\{\bar q_j,\bar p_k\}=\delta_{jk},
\]
并且由泊松括号的反对称性也有
\[
\{\bar p_j,\bar q_k\}=-\delta_{jk}.
\]
这正是 \((2.7.18)\) 式所表达的正则泊松括号关系。它说明一个变量变换若保持哈密顿正则方程的形式不变，则新变量之间必须满足与旧正则变量相同的基本泊松括号关系。
\end{solution}

\begin{solution}[正则变换例子]
\problem{证明到转动参考系的变换
\[
\bar{x}=x\cos\theta-y\sin\theta,
\]
\[
\bar{y}=x\sin\theta+y\cos\theta,
\]
\[
\bar{p}_x=p_x\cos\theta-p_y\sin\theta,
\]
\[
\bar{p}_y=p_x\sin\theta+p_y\cos\theta
\]
是一个正则变换。}
\answer
要证明该变换是正则变换，只需验证新的变量 \((\bar{x},\bar{y},\bar{p}_x,\bar{p}_y)\) 满足与原变量 \((x,y,p_x,p_y)\) 相同的基本泊松括号关系，即
\[
\{\bar{x},\bar{y}\}=0,\qquad \{\bar{p}_x,\bar{p}_y\}=0,
\]
\[
\{\bar{x},\bar{p}_x\}=1,\qquad \{\bar{y},\bar{p}_y\}=1,
\]
\[
\{\bar{x},\bar{p}_y\}=0,\qquad \{\bar{y},\bar{p}_x\}=0.
\]
在二维相空间中，泊松括号定义为
\[
\{f,g\}
=
\frac{\partial f}{\partial x}\frac{\partial g}{\partial p_x}
-\frac{\partial f}{\partial p_x}\frac{\partial g}{\partial x}
+
\frac{\partial f}{\partial y}\frac{\partial g}{\partial p_y}
-\frac{\partial f}{\partial p_y}\frac{\partial g}{\partial y}.
\]
记 \(c=\cos\theta,\ s=\sin\theta\)，则变换可写为
\[
\bar{x}=cx-sy,\qquad \bar{y}=sx+cy,\qquad
\bar{p}_x=cp_x-sp_y,\qquad \bar{p}_y=sp_x+cp_y.
\]
先计算坐标之间的泊松括号。由于 \(\bar{x}\) 和 \(\bar{y}\) 都只依赖 \(x,y\)，不依赖 \(p_x,p_y\)，所以
\[
\{\bar{x},\bar{y}\}=0.
\]
同理，\(\bar{p}_x\) 和 \(\bar{p}_y\) 都只依赖 \(p_x,p_y\)，不依赖 \(x,y\)，所以
\[
\{\bar{p}_x,\bar{p}_y\}=0.
\]
下面计算坐标与对应动量的泊松括号。由
\[
\frac{\partial\bar{x}}{\partial x}=c,\qquad
\frac{\partial\bar{x}}{\partial y}=-s,\qquad
\frac{\partial\bar{p}_x}{\partial p_x}=c,\qquad
\frac{\partial\bar{p}_x}{\partial p_y}=-s,
\]
且 \(\bar{x}\) 不依赖动量、\(\bar{p}_x\) 不依赖坐标，因此
\[
\{\bar{x},\bar{p}_x\}
=
\frac{\partial\bar{x}}{\partial x}\frac{\partial\bar{p}_x}{\partial p_x}
+
\frac{\partial\bar{x}}{\partial y}\frac{\partial\bar{p}_x}{\partial p_y}
=
c^2+s^2=1.
\]
类似地，
\[
\frac{\partial\bar{y}}{\partial x}=s,\qquad
\frac{\partial\bar{y}}{\partial y}=c,\qquad
\frac{\partial\bar{p}_y}{\partial p_x}=s,\qquad
\frac{\partial\bar{p}_y}{\partial p_y}=c,
\]
所以
\[
\{\bar{y},\bar{p}_y\}
=
\frac{\partial\bar{y}}{\partial x}\frac{\partial\bar{p}_y}{\partial p_x}
+
\frac{\partial\bar{y}}{\partial y}\frac{\partial\bar{p}_y}{\partial p_y}
=
s^2+c^2=1.
\]
再计算交叉泊松括号。由于
\[
\frac{\partial\bar{p}_y}{\partial p_x}=s,\qquad
\frac{\partial\bar{p}_y}{\partial p_y}=c,
\]
所以
\[
\{\bar{x},\bar{p}_y\}
=
\frac{\partial\bar{x}}{\partial x}\frac{\partial\bar{p}_y}{\partial p_x}
+
\frac{\partial\bar{x}}{\partial y}\frac{\partial\bar{p}_y}{\partial p_y}
=
cs-sc=0.
\]
同理，
\[
\{\bar{y},\bar{p}_x\}
=
\frac{\partial\bar{y}}{\partial x}\frac{\partial\bar{p}_x}{\partial p_x}
+
\frac{\partial\bar{y}}{\partial y}\frac{\partial\bar{p}_x}{\partial p_y}
=
sc-cs=0.
\]
因此新变量满足
\[
\{\bar{x},\bar{y}\}=0,\qquad
\{\bar{p}_x,\bar{p}_y\}=0,
\]
\[
\{\bar{x},\bar{p}_x\}=1,\qquad
\{\bar{y},\bar{p}_y\}=1,
\]
\[
\{\bar{x},\bar{p}_y\}=0,\qquad
\{\bar{y},\bar{p}_x\}=0.
\]
这正是正则变量所满足的基本泊松括号关系，因此该转动变换是一个正则变换。
\end{solution}


\begin{solution}[正则变换例子]
\problem{证明极坐标的变量 \(\rho=(x^2+y^2)^{1/2}\)，\(\phi=\arctan(y/x)\)，
\[
p_\rho=\hat{\mathbf e}_\rho\cdot \mathbf p=\frac{xp_x+yp_y}{(x^2+y^2)^{1/2}},
\qquad
p_\phi=xp_y-yp_x\;(=l_z)
\]
都是正则变量。这里 \(\hat{\mathbf e}_\rho\) 是径向的单位矢量。}
\answer
要证明 \((\rho,\phi,p_\rho,p_\phi)\) 是正则变量，只需验证它们满足基本泊松括号关系
\[
\{\rho,\phi\}=0,\qquad \{p_\rho,p_\phi\}=0,\qquad
\{\rho,p_\rho\}=1,\qquad \{\phi,p_\phi\}=1,\qquad
\{\rho,p_\phi\}=0,\qquad \{\phi,p_\rho\}=0.
\]
二维相空间中的泊松括号定义为
\[
\{f,g\}=\frac{\partial f}{\partial x}\frac{\partial g}{\partial p_x}-\frac{\partial f}{\partial p_x}\frac{\partial g}{\partial x}+\frac{\partial f}{\partial y}\frac{\partial g}{\partial p_y}-\frac{\partial f}{\partial p_y}\frac{\partial g}{\partial y}.
\]
先列出所需偏导数。由
\[
\rho=(x^2+y^2)^{1/2},\qquad \phi=\arctan(y/x)
\]
可得
\[
\frac{\partial \rho}{\partial x}=\frac{x}{\rho},\qquad
\frac{\partial \rho}{\partial y}=\frac{y}{\rho},\qquad
\frac{\partial \phi}{\partial x}=-\frac{y}{\rho^2},\qquad
\frac{\partial \phi}{\partial y}=\frac{x}{\rho^2}.
\]
同时
\[
p_\rho=\frac{xp_x+yp_y}{\rho},\qquad p_\phi=xp_y-yp_x.
\]
由于 \(\rho\) 和 \(\phi\) 都不依赖动量，所以
\[
\{\rho,\phi\}=0.
\]
下面计算 \(\{\rho,p_\rho\}\)。因为
\[
\frac{\partial p_\rho}{\partial p_x}=\frac{x}{\rho},\qquad
\frac{\partial p_\rho}{\partial p_y}=\frac{y}{\rho},
\]
所以
\[
\{\rho,p_\rho\}
=
\frac{\partial \rho}{\partial x}\frac{\partial p_\rho}{\partial p_x}
+
\frac{\partial \rho}{\partial y}\frac{\partial p_\rho}{\partial p_y}
=
\frac{x^2}{\rho^2}+\frac{y^2}{\rho^2}=1.
\]
再计算 \(\{\phi,p_\phi\}\)。因为
\[
\frac{\partial p_\phi}{\partial p_x}=-y,\qquad
\frac{\partial p_\phi}{\partial p_y}=x,
\]
所以
\[
\{\phi,p_\phi\}
=
\frac{\partial \phi}{\partial x}\frac{\partial p_\phi}{\partial p_x}
+
\frac{\partial \phi}{\partial y}\frac{\partial p_\phi}{\partial p_y}
=
\left(-\frac{y}{\rho^2}\right)(-y)
+
\frac{x}{\rho^2}x
=
\frac{x^2+y^2}{\rho^2}=1.
\]
接着计算交叉括号。首先
\[
\{\rho,p_\phi\}
=
\frac{\partial \rho}{\partial x}\frac{\partial p_\phi}{\partial p_x}
+
\frac{\partial \rho}{\partial y}\frac{\partial p_\phi}{\partial p_y}
=
\frac{x}{\rho}(-y)+\frac{y}{\rho}x=0.
\]
其次
\[
\{\phi,p_\rho\}
=
\frac{\partial \phi}{\partial x}\frac{\partial p_\rho}{\partial p_x}
+
\frac{\partial \phi}{\partial y}\frac{\partial p_\rho}{\partial p_y}
=
-\frac{y}{\rho^2}\frac{x}{\rho}
+
\frac{x}{\rho^2}\frac{y}{\rho}
=0.
\]
最后验证 \(\{p_\rho,p_\phi\}=0\)。先计算
\[
\frac{\partial p_\rho}{\partial x}
=
\frac{p_x}{\rho}-\frac{x(xp_x+yp_y)}{\rho^3},
\qquad
\frac{\partial p_\rho}{\partial y}
=
\frac{p_y}{\rho}-\frac{y(xp_x+yp_y)}{\rho^3},
\]
以及
\[
\frac{\partial p_\phi}{\partial x}=p_y,\qquad
\frac{\partial p_\phi}{\partial y}=-p_x,\qquad
\frac{\partial p_\rho}{\partial p_x}=\frac{x}{\rho},\qquad
\frac{\partial p_\rho}{\partial p_y}=\frac{y}{\rho}.
\]
于是
\[
\begin{aligned}
\{p_\rho,p_\phi\}
&=
\frac{\partial p_\rho}{\partial x}\frac{\partial p_\phi}{\partial p_x}
-\frac{\partial p_\rho}{\partial p_x}\frac{\partial p_\phi}{\partial x}
+\frac{\partial p_\rho}{\partial y}\frac{\partial p_\phi}{\partial p_y}
-\frac{\partial p_\rho}{\partial p_y}\frac{\partial p_\phi}{\partial y}\\
&=
\left(\frac{p_x}{\rho}-\frac{x(xp_x+yp_y)}{\rho^3}\right)(-y)
-\frac{x}{\rho}p_y
+\left(\frac{p_y}{\rho}-\frac{y(xp_x+yp_y)}{\rho^3}\right)x
-\frac{y}{\rho}(-p_x)\\
&=
-\frac{yp_x}{\rho}+\frac{xy(xp_x+yp_y)}{\rho^3}
-\frac{xp_y}{\rho}
+\frac{xp_y}{\rho}-\frac{xy(xp_x+yp_y)}{\rho^3}
+\frac{yp_x}{\rho}
=0.
\end{aligned}
\]
综上，
\[
\{\rho,\phi\}=0,\qquad \{p_\rho,p_\phi\}=0,\qquad
\{\rho,p_\rho\}=1,\qquad \{\phi,p_\phi\}=1,
\]
\[
\{\rho,p_\phi\}=0,\qquad \{\phi,p_\rho\}=0.
\]
这些正是正则变量所要求的基本泊松括号关系，因此 \((\rho,\phi,p_\rho,p_\phi)\) 构成一组正则变量。
\end{solution}



\begin{solution}[正则变换例子]
\problem{证明从变量 \(\mathbf r_1,\mathbf r_2,\mathbf p_1,\mathbf p_2\) 变换到变量 \(\mathbf r_{\rm CM},\mathbf p_{\rm CM},\mathbf r,\mathbf p\) 是一个正则变换。（见习题 2.5.4。）}
\answer
设两粒子的质量分别为 \(m_1,m_2\)，总质量为
\[
M=m_1+m_2,
\]
约化质量为
\[
\mu=\frac{m_1m_2}{m_1+m_2}.
\]
定义质心坐标、相对坐标、总动量和相对动量为
\[
\mathbf r_{\rm CM}=\frac{m_1\mathbf r_1+m_2\mathbf r_2}{M},\qquad
\mathbf r=\mathbf r_1-\mathbf r_2,
\]
\[
\mathbf p_{\rm CM}=\mathbf p_1+\mathbf p_2,\qquad
\mathbf p=\mu(\dot{\mathbf r}_1-\dot{\mathbf r}_2).
\]
由于 \(\mathbf p_1=m_1\dot{\mathbf r}_1\)，\(\mathbf p_2=m_2\dot{\mathbf r}_2\)，所以
\[
\mathbf p
=
\frac{m_1m_2}{M}
\left(
\frac{\mathbf p_1}{m_1}
-
\frac{\mathbf p_2}{m_2}
\right)
=
\frac{m_2\mathbf p_1-m_1\mathbf p_2}{M}.
\]
要证明该变换是正则变换，只需验证新变量满足基本泊松括号关系。记各矢量分量为 \(r_{1\alpha},r_{2\alpha},p_{1\alpha},p_{2\alpha}\)，其中 \(\alpha=x,y,z\)。原变量满足
\[
\{r_{a\alpha},p_{b\beta}\}=\delta_{ab}\delta_{\alpha\beta},\qquad
\{r_{a\alpha},r_{b\beta}\}=0,\qquad
\{p_{a\alpha},p_{b\beta}\}=0,
\]
其中 \(a,b=1,2\)。下面逐一计算新变量的泊松括号。

首先，由于 \(\mathbf r_{\rm CM}\) 和 \(\mathbf r\) 都只依赖坐标而不依赖动量，所以
\[
\{r_\alpha,r_{{\rm CM},\beta}\}=0.
\]
同理，\(\mathbf p\) 和 \(\mathbf p_{\rm CM}\) 都只依赖动量而不依赖坐标，所以
\[
\{p_\alpha,p_{{\rm CM},\beta}\}=0.
\]
接着计算 \(\{r_\alpha,p_\beta\}\)。由
\[
r_\alpha=r_{1\alpha}-r_{2\alpha},\qquad
p_\beta=\frac{m_2p_{1\beta}-m_1p_{2\beta}}{M},
\]
得
\[
\begin{aligned}
\{r_\alpha,p_\beta\}
&=
\left\{r_{1\alpha}-r_{2\alpha},
\frac{m_2p_{1\beta}-m_1p_{2\beta}}{M}
\right\} \\
&=
\frac{m_2}{M}\{r_{1\alpha},p_{1\beta}\}
-\frac{m_1}{M}\{r_{1\alpha},p_{2\beta}\}
-\frac{m_2}{M}\{r_{2\alpha},p_{1\beta}\}
+\frac{m_1}{M}\{r_{2\alpha},p_{2\beta}\} \\
&=
\frac{m_2}{M}\delta_{\alpha\beta}
+\frac{m_1}{M}\delta_{\alpha\beta}
=
\delta_{\alpha\beta}.
\end{aligned}
\]
因此
\[
\{r_\alpha,p_\beta\}=\delta_{\alpha\beta}.
\]
再计算 \(\{r_\alpha,p_{{\rm CM},\beta}\}\)。由
\[
p_{{\rm CM},\beta}=p_{1\beta}+p_{2\beta},
\]
有
\[
\begin{aligned}
\{r_\alpha,p_{{\rm CM},\beta}\}
&=
\{r_{1\alpha}-r_{2\alpha},p_{1\beta}+p_{2\beta}\}\\
&=
\{r_{1\alpha},p_{1\beta}\}
+\{r_{1\alpha},p_{2\beta}\}
-\{r_{2\alpha},p_{1\beta}\}
-\{r_{2\alpha},p_{2\beta}\}\\
&=
\delta_{\alpha\beta}-\delta_{\alpha\beta}=0.
\end{aligned}
\]
所以
\[
\{r_\alpha,p_{{\rm CM},\beta}\}=0.
\]
然后计算 \(\{r_{{\rm CM},\alpha},p_\beta\}\)。由
\[
r_{{\rm CM},\alpha}=\frac{m_1r_{1\alpha}+m_2r_{2\alpha}}{M},\qquad
p_\beta=\frac{m_2p_{1\beta}-m_1p_{2\beta}}{M},
\]
得
\[
\begin{aligned}
\{r_{{\rm CM},\alpha},p_\beta\}
&=
\left\{
\frac{m_1r_{1\alpha}+m_2r_{2\alpha}}{M},
\frac{m_2p_{1\beta}-m_1p_{2\beta}}{M}
\right\}\\
&=
\frac{m_1m_2}{M^2}\{r_{1\alpha},p_{1\beta}\}
-\frac{m_1^2}{M^2}\{r_{1\alpha},p_{2\beta}\}
+\frac{m_2^2}{M^2}\{r_{2\alpha},p_{1\beta}\}
-\frac{m_1m_2}{M^2}\{r_{2\alpha},p_{2\beta}\}\\
&=
\frac{m_1m_2}{M^2}\delta_{\alpha\beta}
-\frac{m_1m_2}{M^2}\delta_{\alpha\beta}
=0.
\end{aligned}
\]
因此
\[
\{r_{{\rm CM},\alpha},p_\beta\}=0.
\]
最后计算 \(\{r_{{\rm CM},\alpha},p_{{\rm CM},\beta}\}\)。有
\[
\begin{aligned}
\{r_{{\rm CM},\alpha},p_{{\rm CM},\beta}\}
&=
\left\{
\frac{m_1r_{1\alpha}+m_2r_{2\alpha}}{M},
p_{1\beta}+p_{2\beta}
\right\}\\
&=
\frac{m_1}{M}\{r_{1\alpha},p_{1\beta}\}
+\frac{m_1}{M}\{r_{1\alpha},p_{2\beta}\}
+\frac{m_2}{M}\{r_{2\alpha},p_{1\beta}\}
+\frac{m_2}{M}\{r_{2\alpha},p_{2\beta}\}\\
&=
\frac{m_1}{M}\delta_{\alpha\beta}
+
\frac{m_2}{M}\delta_{\alpha\beta}
=
\delta_{\alpha\beta}.
\end{aligned}
\]
所以
\[
\{r_{{\rm CM},\alpha},p_{{\rm CM},\beta}\}=\delta_{\alpha\beta}.
\]
综上，新变量满足
\[
\{r_\alpha,p_\beta\}=\delta_{\alpha\beta},\qquad
\{r_{{\rm CM},\alpha},p_{{\rm CM},\beta}\}=\delta_{\alpha\beta},
\]
并且交叉括号为
\[
\{r_\alpha,p_{{\rm CM},\beta}\}=0,\qquad
\{r_{{\rm CM},\alpha},p_\beta\}=0.
\]
同时，新坐标之间以及新动量之间的泊松括号也均为零：
\[
\{r_\alpha,r_\beta\}=0,\qquad
\{r_{{\rm CM},\alpha},r_{{\rm CM},\beta}\}=0,\qquad
\{r_\alpha,r_{{\rm CM},\beta}\}=0,
\]
\[
\{p_\alpha,p_\beta\}=0,\qquad
\{p_{{\rm CM},\alpha},p_{{\rm CM},\beta}\}=0,\qquad
\{p_\alpha,p_{{\rm CM},\beta}\}=0.
\]
因此 \((\mathbf r,\mathbf p)\) 和 \((\mathbf r_{\rm CM},\mathbf p_{\rm CM})\) 分别构成两组正则共轭变量，且两组变量彼此泊松对易。故从
\(
(\mathbf r_1,\mathbf r_2,\mathbf p_1,\mathbf p_2)
\)
到
\(
(\mathbf r_{\rm CM},\mathbf p_{\rm CM},\mathbf r,\mathbf p)
\)
的变换是一个正则变换。
\end{solution}



\begin{solution}[正则变换例子]
\problem{证明
\[
\bar q=\ln(q^{-1}\sin p),
\]
\[
\bar p=q\cot p
\]
是一个正则变换。}
\answer
由于这是单自由度系统，原变量为 \((q,p)\)，新变量为 \((\bar q,\bar p)\)。在单自由度情形下，泊松括号定义为
\[
\{f,g\}
=
\frac{\partial f}{\partial q}\frac{\partial g}{\partial p}
-
\frac{\partial f}{\partial p}\frac{\partial g}{\partial q}.
\]
要证明变换是正则变换，只需验证
\[
\{\bar q,\bar p\}=1.
\]
因为单自由度下任意函数与自身的泊松括号恒为零，所以自动有
\[
\{\bar q,\bar q\}=0,\qquad \{\bar p,\bar p\}=0.
\]
现在计算 \(\{\bar q,\bar p\}\)。由
\[
\bar q=\ln(q^{-1}\sin p)=\ln(\sin p)-\ln q
\]
可得
\[
\frac{\partial \bar q}{\partial q}=-\frac{1}{q},\qquad
\frac{\partial \bar q}{\partial p}=\cot p.
\]
又因为
\[
\bar p=q\cot p,
\]
所以
\[
\frac{\partial \bar p}{\partial q}=\cot p,\qquad
\frac{\partial \bar p}{\partial p}=-q\csc^2 p.
\]
代入泊松括号定义，得到
\[
\begin{aligned}
\{\bar q,\bar p\}
&=
\frac{\partial \bar q}{\partial q}\frac{\partial \bar p}{\partial p}
-
\frac{\partial \bar q}{\partial p}\frac{\partial \bar p}{\partial q}\\
&=
\left(-\frac{1}{q}\right)(-q\csc^2 p)-(\cot p)(\cot p)\\
&=
\csc^2 p-\cot^2 p.
\end{aligned}
\]
利用恒等式
\[
\csc^2 p-\cot^2 p=1,
\]
可得
\[
\{\bar q,\bar p\}=1.
\]
因此新变量满足基本正则泊松括号关系
\[
\{\bar q,\bar q\}=0,\qquad
\{\bar p,\bar p\}=0,\qquad
\{\bar q,\bar p\}=1.
\]
所以该变换是一个正则变换。
\end{solution}




\begin{solution}[正则变换中的动量变换]
\problem{我们想在这里推导出 \((2.7.9)\) 式，它给出了组态空间中坐标
\[
q_i\to \bar q_i(q_1,\cdots,q_n)
\]
的变换下，动量的变换。

(1) 如果我们把上面的方程反解得到 \(q=q(\bar q)\)，我们就可以导出下面的 \((2.7.7)\) 式的对应式：
\[
\dot q_i=\sum_j\left(\frac{\partial q_i}{\partial \bar q_j}\right)\dot{\bar q}_j.
\]

(2) 从上式证明
\[
\left(\frac{\partial \dot q_i}{\partial \dot{\bar q}_j}\right)_{\bar q}
=
\frac{\partial q_i}{\partial \bar q_j}.
\]

(3) 现在计算
\[
\bar p_i=
\left[\frac{\partial \mathcal L(\bar q,\dot{\bar q})}{\partial \dot{\bar q}_i}\right]_{\bar q}
=
\left[\frac{\partial \mathcal L(q,\dot q)}{\partial \dot{\bar q}_i}\right]_{\bar q}.
\]
利用链式规则，以及 \(q=q(\bar q)\) 而不是 \(q(\bar q,\dot{\bar q})\)，推导 \((2.7.9)\) 式。

(4) 通过计算 \((2.7.18)\) 式中的泊松括号，证明点变换是正则变换。}
\answer
设组态空间中的点变换为
\[
\bar q_i=\bar q_i(q_1,\cdots,q_n),
\]
并假设该变换可逆，因此可以反解为
\[
q_i=q_i(\bar q_1,\cdots,\bar q_n).
\]
因为 \(q_i\) 只依赖新的坐标 \(\bar q_j\)，不依赖新的速度 \(\dot{\bar q}_j\)，所以对时间求导时由链式法则得到
\[
\dot q_i=\frac{d q_i(\bar q)}{dt}
=
\sum_j\frac{\partial q_i}{\partial \bar q_j}\dot{\bar q}_j.
\]
这就是 \((2.7.7)\) 式在反变换 \(q=q(\bar q)\) 下的对应形式。

下面证明
\[
\left(\frac{\partial \dot q_i}{\partial \dot{\bar q}_j}\right)_{\bar q}
=
\frac{\partial q_i}{\partial \bar q_j}.
\]
由刚才得到的速度变换式
\[
\dot q_i=\sum_k\frac{\partial q_i}{\partial \bar q_k}\dot{\bar q}_k
\]
可知，在对 \(\dot{\bar q}_j\) 求偏导时，保持 \(\bar q\) 不变。由于 \(\partial q_i/\partial \bar q_k\) 只是 \(\bar q\) 的函数，在该偏导中应视为常数，因此
\[
\left(\frac{\partial \dot q_i}{\partial \dot{\bar q}_j}\right)_{\bar q}
=
\sum_k\frac{\partial q_i}{\partial \bar q_k}
\frac{\partial \dot{\bar q}_k}{\partial \dot{\bar q}_j}
=
\sum_k\frac{\partial q_i}{\partial \bar q_k}\delta_{kj}
=
\frac{\partial q_i}{\partial \bar q_j}.
\]
接下来推导动量变换公式。旧变量下的共轭动量为
\[
p_i=\frac{\partial \mathcal L(q,\dot q)}{\partial \dot q_i},
\]
新变量下的共轭动量为
\[
\bar p_j=\frac{\partial \mathcal L(\bar q,\dot{\bar q})}{\partial \dot{\bar q}_j}.
\]
由于只是点变换，\(q=q(\bar q)\)，并且 \(\dot q=\dot q(\bar q,\dot{\bar q})\)，所以
\[
\mathcal L(\bar q,\dot{\bar q})
=
\mathcal L\bigl(q(\bar q),\dot q(\bar q,\dot{\bar q})\bigr).
\]
对 \(\dot{\bar q}_j\) 求偏导，并保持 \(\bar q\) 不变。因为 \(q=q(\bar q)\) 不依赖 \(\dot{\bar q}\)，所以链式法则中只有 \(\dot q_i\) 项贡献：
\[
\bar p_j
=
\left(\frac{\partial \mathcal L}{\partial \dot{\bar q}_j}\right)_{\bar q}
=
\sum_i
\frac{\partial \mathcal L}{\partial \dot q_i}
\left(\frac{\partial \dot q_i}{\partial \dot{\bar q}_j}\right)_{\bar q}.
\]
利用
\[
\frac{\partial \mathcal L}{\partial \dot q_i}=p_i,
\qquad
\left(\frac{\partial \dot q_i}{\partial \dot{\bar q}_j}\right)_{\bar q}
=
\frac{\partial q_i}{\partial \bar q_j},
\]
得到
\[
\bar p_j
=
\sum_i p_i\frac{\partial q_i}{\partial \bar q_j}.
\]
即
\[
\boxed{\bar p_j=\sum_i\frac{\partial q_i}{\partial \bar q_j}p_i.}
\]
这就是点变换下动量的变换规律，也就是 \((2.7.9)\) 式。它表明动量按照坐标变换雅可比矩阵的逆变换进行变换。

最后证明该点变换是正则变换。需要验证
\[
\{\bar q_i,\bar q_j\}=0,\qquad
\{\bar p_i,\bar p_j\}=0,\qquad
\{\bar q_i,\bar p_j\}=\delta_{ij}.
\]
首先，因为 \(\bar q_i\) 和 \(\bar q_j\) 都只依赖旧坐标 \(q\)，不依赖旧动量 \(p\)，所以
\[
\{\bar q_i,\bar q_j\}=0.
\]
下面计算 \(\{\bar q_i,\bar p_j\}\)。由
\[
\bar p_j=\sum_k\frac{\partial q_k}{\partial \bar q_j}p_k
\]
可得
\[
\{\bar q_i,\bar p_j\}
=
\sum_k
\frac{\partial q_k}{\partial \bar q_j}
\{\bar q_i,p_k\},
\]
其中 \(\partial q_k/\partial \bar q_j\) 只依赖坐标，和 \(\bar q_i\) 的泊松括号为零。又因为
\[
\{\bar q_i,p_k\}
=
\frac{\partial \bar q_i}{\partial q_k},
\]
所以
\[
\{\bar q_i,\bar p_j\}
=
\sum_k
\frac{\partial \bar q_i}{\partial q_k}
\frac{\partial q_k}{\partial \bar q_j}.
\]
根据反函数的链式法则，
\[
\sum_k
\frac{\partial \bar q_i}{\partial q_k}
\frac{\partial q_k}{\partial \bar q_j}
=
\frac{\partial \bar q_i}{\partial \bar q_j}
=
\delta_{ij}.
\]
因此
\[
\{\bar q_i,\bar p_j\}=\delta_{ij}.
\]
最后，由泊松括号的反对称性可得
\[
\{\bar p_i,\bar q_j\}=-\delta_{ij}.
\]
再证明 \(\{\bar p_i,\bar p_j\}=0\)。这可以直接由正则结构推出，也可以显式计算。写成微分形式更简洁：由
\[
\bar p_j=\sum_i p_i\frac{\partial q_i}{\partial \bar q_j}
\]
可得
\[
\sum_j\bar p_j\,d\bar q_j
=
\sum_j\sum_i p_i\frac{\partial q_i}{\partial \bar q_j}d\bar q_j
=
\sum_i p_i\,dq_i.
\]
因此正则一形式保持不变：
\[
\sum_i p_i\,dq_i=\sum_j\bar p_j\,d\bar q_j.
\]
对两边取外微分，得到
\[
\sum_i dp_i\wedge dq_i
=
\sum_j d\bar p_j\wedge d\bar q_j,
\]
即辛结构保持不变。因此点变换保持基本泊松括号关系，特别有
\[
\{\bar p_i,\bar p_j\}=0.
\]
综上，
\[
\{\bar q_i,\bar q_j\}=0,\qquad
\{\bar p_i,\bar p_j\}=0,\qquad
\{\bar q_i,\bar p_j\}=\delta_{ij}.
\]
所以由组态空间点变换
\[
q_i\to \bar q_i(q_1,\cdots,q_n)
\]
以及动量变换
\[
\bar p_j=\sum_i\frac{\partial q_i}{\partial \bar q_j}p_i
\]
所构成的相空间变换是正则变换。
\end{solution}


\begin{solution}[泊松括号不变性]
\problem{通过直接计算验证方程 \((2.7.19)\)。使用链式规则从对 \(q,p\) 的导数变成对 \(\bar q,\bar p\) 的导数。收集表示后者的 PB 的各项。}
\answer
要证明的是，在正则变换
\[
(q,p)\longrightarrow (\bar q,\bar p)
\]
下，泊松括号的形式保持不变，即
\[
\{\omega,\sigma\}_{q,p}=\{\omega,\sigma\}_{\bar q,\bar p}.
\]
这里下标表示泊松括号是用哪一组正则变量计算的。为简化记号，先考虑单自由度情形。按照定义，
\[
\{\omega,\sigma\}_{q,p}
=
\frac{\partial \omega}{\partial q}\frac{\partial \sigma}{\partial p}
-
\frac{\partial \omega}{\partial p}\frac{\partial \sigma}{\partial q}.
\]
由于 \(\omega\) 和 \(\sigma\) 可以看成新变量 \(\bar q,\bar p\) 的函数，即
\[
\omega=\omega(\bar q,\bar p),\qquad \sigma=\sigma(\bar q,\bar p),
\]
而 \(\bar q,\bar p\) 又是 \(q,p\) 的函数，因此由链式法则有
\[
\frac{\partial \omega}{\partial q}
=
\frac{\partial \omega}{\partial \bar q}\frac{\partial \bar q}{\partial q}
+
\frac{\partial \omega}{\partial \bar p}\frac{\partial \bar p}{\partial q},
\qquad
\frac{\partial \omega}{\partial p}
=
\frac{\partial \omega}{\partial \bar q}\frac{\partial \bar q}{\partial p}
+
\frac{\partial \omega}{\partial \bar p}\frac{\partial \bar p}{\partial p},
\]
同理
\[
\frac{\partial \sigma}{\partial q}
=
\frac{\partial \sigma}{\partial \bar q}\frac{\partial \bar q}{\partial q}
+
\frac{\partial \sigma}{\partial \bar p}\frac{\partial \bar p}{\partial q},
\qquad
\frac{\partial \sigma}{\partial p}
=
\frac{\partial \sigma}{\partial \bar q}\frac{\partial \bar q}{\partial p}
+
\frac{\partial \sigma}{\partial \bar p}\frac{\partial \bar p}{\partial p}.
\]
将这些表达式代入 \(\{\omega,\sigma\}_{q,p}\)，得到
\[
\begin{aligned}
\{\omega,\sigma\}_{q,p}
={}&
\left(
\frac{\partial \omega}{\partial \bar q}\frac{\partial \bar q}{\partial q}
+
\frac{\partial \omega}{\partial \bar p}\frac{\partial \bar p}{\partial q}
\right)
\left(
\frac{\partial \sigma}{\partial \bar q}\frac{\partial \bar q}{\partial p}
+
\frac{\partial \sigma}{\partial \bar p}\frac{\partial \bar p}{\partial p}
\right)\\
&-
\left(
\frac{\partial \omega}{\partial \bar q}\frac{\partial \bar q}{\partial p}
+
\frac{\partial \omega}{\partial \bar p}\frac{\partial \bar p}{\partial p}
\right)
\left(
\frac{\partial \sigma}{\partial \bar q}\frac{\partial \bar q}{\partial q}
+
\frac{\partial \sigma}{\partial \bar p}\frac{\partial \bar p}{\partial q}
\right).
\end{aligned}
\]
展开后，含有
\[
\frac{\partial \omega}{\partial \bar q}
\frac{\partial \sigma}{\partial \bar q}
\]
的项相互抵消，含有
\[
\frac{\partial \omega}{\partial \bar p}
\frac{\partial \sigma}{\partial \bar p}
\]
的项也相互抵消，剩下
\[
\begin{aligned}
\{\omega,\sigma\}_{q,p}
={}&
\frac{\partial \omega}{\partial \bar q}
\frac{\partial \sigma}{\partial \bar p}
\left(
\frac{\partial \bar q}{\partial q}\frac{\partial \bar p}{\partial p}
-
\frac{\partial \bar q}{\partial p}\frac{\partial \bar p}{\partial q}
\right)\\
&+
\frac{\partial \omega}{\partial \bar p}
\frac{\partial \sigma}{\partial \bar q}
\left(
\frac{\partial \bar p}{\partial q}\frac{\partial \bar q}{\partial p}
-
\frac{\partial \bar p}{\partial p}\frac{\partial \bar q}{\partial q}
\right).
\end{aligned}
\]
括号中的第一项正是
\[
\{\bar q,\bar p\}_{q,p}
=
\frac{\partial \bar q}{\partial q}\frac{\partial \bar p}{\partial p}
-
\frac{\partial \bar q}{\partial p}\frac{\partial \bar p}{\partial q},
\]
而第二项是
\[
\{\bar p,\bar q\}_{q,p}
=
\frac{\partial \bar p}{\partial q}\frac{\partial \bar q}{\partial p}
-
\frac{\partial \bar p}{\partial p}\frac{\partial \bar q}{\partial q}.
\]
若变换是正则变换，则
\[
\{\bar q,\bar p\}_{q,p}=1,\qquad
\{\bar p,\bar q\}_{q,p}=-1.
\]
因此
\[
\begin{aligned}
\{\omega,\sigma\}_{q,p}
&=
\frac{\partial \omega}{\partial \bar q}
\frac{\partial \sigma}{\partial \bar p}
-
\frac{\partial \omega}{\partial \bar p}
\frac{\partial \sigma}{\partial \bar q}\\
&=
\{\omega,\sigma\}_{\bar q,\bar p}.
\end{aligned}
\]
所以
\[
\boxed{\{\omega,\sigma\}_{q,p}=\{\omega,\sigma\}_{\bar q,\bar p}.}
\]
这说明在正则变换下，泊松括号的形式保持不变。对于多自由度系统，证明完全类似，只需把单自由度的泊松括号推广为
\[
\{f,g\}
=
\sum_i
\left(
\frac{\partial f}{\partial q_i}\frac{\partial g}{\partial p_i}
-
\frac{\partial f}{\partial p_i}\frac{\partial g}{\partial q_i}
\right),
\]
并利用正则关系
\[
\{\bar q_i,\bar q_j\}=0,\qquad
\{\bar p_i,\bar p_j\}=0,\qquad
\{\bar q_i,\bar p_j\}=\delta_{ij}.
\]
由此同样得到泊松括号在正则变量变换下不变。
\end{solution}